\chapter{Introdución}
\label{chap:introducion}

\lettrine{P}{rimer} capítulo de la memoria ... Comienza con el marco o contexto general en el que se sitúa la problemática a resolver por este proyecto. 

{ \color{red} Aspectos a considerar *en general*
    \begin{itemize}
      \item Redacción en impersonal preferentemente o primer persona del plural (nunca en primer persona del singular).
      \item Muy importante, todas las figuras o tablas tienen que estar referenciadas desde el texto, así como comentadas.
      \item Referencias a bibliografía / citas desde texto y cuando se introduzca un nuevo concepto (similar con acrónimos, la primera vez que aparezca un acrónimo, se indicará entre paréntesis la versión completa del mismo; el resto de veces ya se puede indicar directamente el acrónimo).
    \end{itemize}
}

\section{Determinación de la situación actual}
\label{sec:actual}

{ \color{red} Aspectos a considerar
    \begin{itemize}
      \item En la sección sección \ref{sec:actual}, página \pageref{sec:actual} se comenta la descripción de la situación actual, cómo se aborda la resolución del problema sin la introducción del sistema a desarrollar en este trabajo. Esta sección termina con la conclusión de la necesidad de una aplicación web para resolver este problema. Para este apartado es posible abordarlo realizando un análisis DAFO (\url{https://es.wikipedia.org/wiki/An%C3%A1lisis_DAFO} ) de la situación actual para un problema, de modo que se identifiquen las posibles mejoras que darán lugar al TFG.
    \end{itemize}
}

\section{Alcance y objetivos}
\label{sec:objetivos}

{ \color{red} Aspectos a considerar
    \begin{itemize}
      \item En la sección \ref{sec:objetivos}, página \pageref{sec:objetivos}, una vez contextualizado el problema y las formas en las que se resuelve actualmente (de forma manual, con ayuda de otras aplicaciones que no resuelven el problema de forma completo), se comentan los objetivos concretos, que no los requisitos completos (eso se deja para el capítulo \pageref{chap:requisitos} ).
    \end{itemize}
}

\subsection{Objetivos principales}

\begin{enumerate}
    \item Primer objetivo del proyecto
    \item Segundo objetivo del proyecto
    \item Los siguientes si los hubiere \ldots
\end{enumerate}


\subsection{Objetivos concretos}

\begin{enumerate}
    \item Primero
    \item Segundo 
    \item Y así sucesivamente \ldots
\end{enumerate}


\section{Estructura de la memoria}
\label{sec:estructura}
Breve  explicación de que se puede encontrar en cada capítulo.

\begin{enumerate}
    \item Este mismo capítulo.
    \item El segundo de los capítulo centra \ldots
\end{enumerate}

